%input{vars.tex}
\usepackage{ifdraft}
\usepackage{luacode}
\usepackage{ifmtarg}
\usepackage{environ}
\usepackage{hyperref}
\usepackage{hyperxmp}
\usepackage{lipsum}
\usepackage{shellesc}
\usepackage{booktabs}
\usepackage{tabularx}

\usepackage{fourier}
\usepackage{heuristica}
\usepackage{array}
%usepackage{enumitem}   
\usepackage{tabularray}   
\UseTblrLibrary{booktabs, varwidth} 

\usepackage{tcolorbox}
\tcbuselibrary{skins}
\usepackage{fontawesome5}
\usepackage{pdftexcmds}
\usepackage{robust-externalize}
\ExplSyntaxOn
\ProvideExpandableDocumentCommand{\lenToCmNoUnit}{ O{cm} m }
  { \dim_to_decimal_in_unit:nn { #2 } { 1 #1 } }
\ExplSyntaxOff
\usepackage{relsize}
\usepackage{algorithm}
\usepackage{verbatimbox}
\usepackage{forest}
\usepackage{tkz-berge}
\if\dataeq{type}{technicalslides}
\PassOptionsToPackage{noend}{algpseudocode}
\usepackage{algpseudocode}
\else
\PassOptionsToPackage{noEnd,indLines=false,italicComments=false}{algpseudocodex}
\usepackage{algpseudocodex}
\fi
\usepackage{float}
\usepackage{needspace}
\usepackage{soul}
\usepackage{setspace}
\usepackage{mathtools}
\usepackage{caption}
\usepackage[dvipsnames, x11names, svgnames]{xcolor}
\usepackage{colortbl}
\usepackage{color,soul}
\usepackage{listings}
\usepackage{minted}
\usepackage{filecontents}
\usepackage{nowidow}
\usepackage{amssymb}
\usepackage{amsmath}
\usepackage{stmaryrd}

% ---------------------------------------------------------------------------
% Source-code references and authoring notices
% ---------------------------------------------------------------------------

\definecolor{BioseqCodeRefIcon}{HTML}{24292F}
\definecolor{BioseqCodeRefTitle}{HTML}{0969DA}
\definecolor{BioseqCodeRefDetail}{HTML}{57606A}

% Repository aliases keep source references concise and make adding another
% upstream a single declaration.  The optional third argument records a
% non-default branch; repositories otherwise target `main`.
\ExplSyntaxOn
\prop_new:N \g__bioseq_code_repositories_prop
\prop_new:N \g__bioseq_code_repository_branches_prop
\seq_new:N \l__bioseq_coderef_parts_seq
\tl_new:N \l__bioseq_coderef_repo_tl
\tl_new:N \l__bioseq_coderef_path_tl
\tl_new:N \l__bioseq_coderef_base_tl
\tl_new:N \l__bioseq_coderef_branch_tl
\tl_new:N \l__bioseq_coderef_url_tl
\tl_new:N \l__bioseq_coderef_description_tl
\tl_new:N \l__bioseq_coderef_location_tl

\NewDocumentCommand{\DeclareCodeRepository}{m m O{main}}
  {
    \prop_gput:Nnn \g__bioseq_code_repositories_prop {#1} {#2}
    \prop_gput:Nnn \g__bioseq_code_repository_branches_prop {#1} {#3}
  }

\DeclareCodeRepository{minimap2}{https://github.com/lh3/minimap2}[master]
\DeclareCodeRepository{g3sa}{https://github.com/sunwookim028/G3SA}
\DeclareCodeRepository{sneakysnake}{https://github.com/CMU-SAFARI/SneakySnake}
\DeclareCodeRepository{blockaligner}{https://github.com/Daniel-Liu-c0deb0t/block-aligner}
\DeclareCodeRepository{block-aligner}{https://github.com/Daniel-Liu-c0deb0t/block-aligner}
\DeclareCodeRepository{block_aligner}{https://github.com/Daniel-Liu-c0deb0t/block-aligner}
\DeclareCodeRepository{postgres}{https://github.com/postgres/postgres}
\DeclareCodeRepository{postgresql}{https://github.com/postgres/postgres}

\cs_generate_variant:Nn \prop_get:NnNTF { NVN }

\cs_new_protected:Npn \__bioseq_coderef_link:n #1
  {
    \href{#1}{%
      \mbox{%
        \textcolor{BioseqCodeRefIcon}{\fontsize{7.5}{7.5}\selectfont\faGithub}%
        \hspace{.3em}%
        \begin{tabular}[c]{@{}l@{}}
          \sffamily\fontsize{6.5}{7}\selectfont\bfseries
          \textcolor{BioseqCodeRefTitle}{\tl_use:N \l__bioseq_coderef_description_tl}
          \\[-.55ex]
          \ttfamily\fontsize{5.5}{6}\selectfont
          \textcolor{BioseqCodeRefDetail}{\tl_use:N \l__bioseq_coderef_location_tl}
        \end{tabular}%
      }%
    }%
  }

\NewDocumentCommand{\coderef}{m m m m g}
  {
    \group_begin:
      \seq_set_split:Nnn \l__bioseq_coderef_parts_seq {/} {#1}
      \seq_pop_left:NN \l__bioseq_coderef_parts_seq \l__bioseq_coderef_repo_tl
      \tl_set:Nx \l__bioseq_coderef_path_tl
        { \seq_use:Nn \l__bioseq_coderef_parts_seq {/} }
      \prop_get:NVNTF
        \g__bioseq_code_repositories_prop
        \l__bioseq_coderef_repo_tl
        \l__bioseq_coderef_base_tl
        { }
        {
          \PackageError{bioseq-common}
            {Unknown~code~repository~alias~'\tl_to_str:V \l__bioseq_coderef_repo_tl'}
            {Declare~it~with~\string\DeclareCodeRepository\{alias\}\{URL\}.}
        }
      \prop_get:NVNTF
        \g__bioseq_code_repository_branches_prop
        \l__bioseq_coderef_repo_tl
        \l__bioseq_coderef_branch_tl
        { }
        { \tl_set:Nn \l__bioseq_coderef_branch_tl {main} }
      \tl_if_empty:NT \l__bioseq_coderef_path_tl
        {
          \PackageError{bioseq-common}
            {A~code~reference~requires~alias/path}
            {Pass~the~repository~alias,~a~slash,~and~the~path~inside~the~repository.}
        }
      \tl_set:Nx \l__bioseq_coderef_url_tl
        {
          \tl_use:N \l__bioseq_coderef_base_tl
          /blob/\tl_use:N \l__bioseq_coderef_branch_tl
          /\tl_use:N \l__bioseq_coderef_path_tl
          \c_hash_str L\tl_to_str:n {#4}
        }
      \tl_set:Nn \l__bioseq_coderef_description_tl {#2}
      \tl_set:Nx \l__bioseq_coderef_location_tl
        { \tl_to_str:n {#3} \c_space_tl \tl_to_str:n {#4} }
      \IfNoValueF{#5}
        { \tl_put_right:Nx \l__bioseq_coderef_location_tl { :\tl_to_str:n {#5} } }
      \unskip\penalty0\hspace*{.6em plus 1fill}
      \exp_args:NV \__bioseq_coderef_link:n \l__bioseq_coderef_url_tl
    \group_end:
  }
\ExplSyntaxOff

\definecolor{BioseqTodoBackground}{HTML}{FFF8C5}
\definecolor{BioseqTodoFrame}{HTML}{9A6700}
\definecolor{BioseqWarningBackground}{HTML}{FFF4CE}
\definecolor{BioseqWarningFrame}{HTML}{9A6700}
\definecolor{BioseqInfoBackground}{HTML}{DDF4FF}
\definecolor{BioseqInfoFrame}{HTML}{0969DA}

% Use tcbox's explicit argument form: body-capturing tcolorbox environments
% cannot be hidden inside a command because their end delimiter is then lost.
\NewDocumentCommand{\todo}{m}{%
  \par\addvspace{.55\baselineskip}%
  \begingroup
  \let\tikzpicture\oldtikzpicture
  \let\endtikzpicture\oldendtikzpicture
    \noindent\tcbox[
    enhanced,nobeforeafter,width=\linewidth,tcbox width=auto limited,
    colback=BioseqTodoBackground,colframe=BioseqTodoFrame,
    colbacktitle=BioseqTodoBackground,coltitle=BioseqTodoFrame,
    boxrule=.55pt,arc=1.5mm,
    left=7pt,right=7pt,top=5pt,bottom=5pt,
    title={\faStickyNote\hspace{.5em}TODO},
    fonttitle=\sffamily\bfseries
  ]{#1}%
  \endgroup
  \par\addvspace{.55\baselineskip}%
}

\NewDocumentCommand{\warn}{m +m}{%
  \par\addvspace{.65\baselineskip}%
  \begingroup
  \let\tikzpicture\oldtikzpicture
  \let\endtikzpicture\oldendtikzpicture
    \noindent\tcbox[capture=minipage,
    enhanced,nobeforeafter,width=\linewidth,tcbox width=forced left,
    colback=BioseqWarningBackground,colframe=BioseqWarningFrame,
    colbacktitle=BioseqWarningBackground,coltitle=BioseqWarningFrame,
    borderline west={2.2pt}{0pt}{BioseqWarningFrame},
    boxrule=.45pt,arc=1.5mm,
    left=8pt,right=8pt,top=6pt,bottom=6pt,
    title={\faExclamationTriangle\hspace{.55em}#1},
    fonttitle=\sffamily\bfseries
  ]{#2}%
  \endgroup
  \par\addvspace{.65\baselineskip}%
}

\NewDocumentCommand{\info}{m m}{%
  \par\addvspace{.65\baselineskip}%
  \begingroup
  \let\tikzpicture\oldtikzpicture
  \let\endtikzpicture\oldendtikzpicture
    \noindent\tcbox[
    enhanced,nobeforeafter,width=\linewidth,tcbox width=auto limited,
    colback=BioseqInfoBackground,colframe=BioseqInfoFrame,
    colbacktitle=BioseqInfoBackground,coltitle=BioseqInfoFrame,
    borderline west={2.2pt}{0pt}{BioseqInfoFrame},
    boxrule=.45pt,arc=1.5mm,
    left=8pt,right=8pt,top=6pt,bottom=6pt,
    title={\faInfoCircle\hspace{.55em}#1},
    fonttitle=\sffamily\bfseries
  ]{#2}%
  \endgroup
  \par\addvspace{.65\baselineskip}%
}

% Author notes are deliberately visible in the PDF.  TeX comments such as
% `% TODO: ...` are discarded by TeX and therefore cannot serve this purpose.

\usepackage{smartdiagram}
\usesmartdiagramlibrary {core.definitions}
\usesmartdiagramlibrary { core.styles }
\usesmartdiagramlibrary { core.commands }
\usepackage{bytefield}

\usepackage{emoji}
\usepackage{fontspec}

\usepackage{ulem}
\usepackage{texshade} 

\directlua{luaotfload.add_fallback
   ("emojifallback",
    {
        "Apple Color Emoji:mode=harf;",
        "Noto Color Emoji:mode=harf;",
    })}

\setnowidow
\setnoclub


\pgfmathsetseed{\number\randomseed}
\usetikzlibrary{graphs, graphdrawing, quotes, tikzmark, positioning, decorations.pathreplacing, decorations.pathmorphing, calc, arrows.meta, shapes.geometric, backgrounds, arrows, automata, petri, shapes, snakes, mindmap}
%\usegdlibrary{layered}
\usegdlibrary{force}

\definecolor{lightblue}{RGB}{160,180,200}
\definecolor{darkblue}{RGB}{110,130,150}

\begin{PlaceholderFromCode}{__EXT_PYPLOT_TEMPLATE_BEFORE__}
  import matplotlib.pyplot as plt
  import matplotlib
  import numpy as np
  import json
  matplotlib.use("pgf")
  matplotlib.rcParams.update({
    "font.family": "serif",
    "font.serif": [],
    "text.usetex": True,
  })
  fig, ax = plt.subplots()
  fig.set_size_inches(__LINEWIDTH__, 0.7*__LINEWIDTH__, forward=True)
\end{PlaceholderFromCode}
\begin{PlaceholderFromCode}{__EXT_PYPLOT_TEMPLATE_AFTER__}
  print(get_filename_from_extension(".pgf"))
  plt.savefig(get_filename_from_extension(".pgf"), bbox_inches="tight")
\end{PlaceholderFromCode}

\robExtConfigure{
  compile in parallel with gnu parallel,
  new preset={pyplot}{
    python,
    % We create a new placeholder (it is simple enough that you don't need to use PlaceholderFromCode)
    add before placeholder no space={__ROBEXT_MAIN_CONTENT__}{__EXT_PYPLOT_TEMPLATE_BEFORE__},
    add to placeholder no space={__ROBEXT_MAIN_CONTENT__}{__EXT_PYPLOT_TEMPLATE_AFTER__},
    custom include command={
      \begin{figure}[H]
        \centering
        \input{\robExtAddCachePathAndName{\robExtFinalHash.pgf}}
      \end{figure}
    },
    set placeholder eval={__LINEWIDTH__}{\lenToCmNoUnit[in]{\linewidth}}
  },
}

\def\mathdefault#1{#1} % Needed in matplotlib 3.8: https://github.com/matplotlib/matplotlib/issues/27907


\definecolor{BioseqAlgorithmCommentBackground}{HTML}{FFF3B0}
\definecolor{BioseqAlgorithmCommentText}{HTML}{6E4C00}
\newcommand{\AlgComment}[1]{%
  \begingroup
  \setlength{\fboxsep}{1.2pt}%
  \colorbox{BioseqAlgorithmCommentBackground}{%
    \strut\normalfont\scriptsize\sffamily
    \color{BioseqAlgorithmCommentText}#1}%
  \endgroup
}
\renewcommand{\algorithmiccomment}[1]{\hfill\AlgComment{#1}}
\newcommand{\CommentH}[1]{\Comment{\textbf{#1}}}

\tikzset{
  fontscale/.style={font=\relsize{#1}},
  algpxIndentLine/.style={draw=gray!65,line width=.35pt}
}

\DeclarePairedDelimiter\ceil{\lceil}{\rceil}
\DeclarePairedDelimiter\floor{\lfloor}{\rfloor}

%% FIXES FOR ALGORITHMS

\usepackage{etoolbox}

\newcommand{\algruledefaultfactor}{.75}
\newcommand{\algstrut}[1][\algruledefaultfactor]{\vrule width 0pt
depth .25\baselineskip height #1\baselineskip\relax}
\newcommand*{\algrule}[1][\algorithmicindent]{%
  \hspace*{.5em}%
  % Slightly overlap adjacent row segments.  This removes hairline seams from
  % PDF antialiasing without changing the algorithm's row height.
  \smash{\raisebox{-.30\baselineskip}{\rule{.4pt}{1.10\baselineskip}}}%
  \algstrut
  \hspace*{\dimexpr#1-.5em}}

\makeatletter
\newcount\ALG@printindent@tempcnta
\def\ALG@printindent@rules{%
    \unskip
    % draw a rule for each indent level
    \ALG@printindent@tempcnta=1
    \loop
    \algrule[\csname ALG@ind@\the\ALG@printindent@tempcnta\endcsname]%
    \advance \ALG@printindent@tempcnta 1
    \ifnum \ALG@printindent@tempcnta<\numexpr\theALG@nested+1\relax% can't do <=, so add one to RHS and use < instead
    \repeat
}
\def\ALG@printindent{%
    \ifnum \theALG@nested>0% is there anything to print
    \ifx\ALG@text\ALG@x@notext% hidden end groups must not create a row
    % do nothing
    \else
    \ALG@printindent@rules
    \fi
    \fi
}%

\patchcmd{\ALG@doentity}{\noindent\hskip\ALG@tlm}{\ALG@printindent}{}{\errmessage{failed to patch}}
% algorithmicx normally emits an empty list item for every hidden EndFor,
% EndIf, and EndWhile. With noEnd enabled those phantom rows create large
% blank separators between otherwise adjacent pseudocode phases.
\patchcmd{\ALG@doentity}{\item[]\nointerlineskip}{\relax}{}{\errmessage{failed to suppress hidden algorithm row}}

\newcommand{\AlgCommentLine}[1]{%
  \Statex
  \ifnum\theALG@nested>0\ALG@printindent@rules\fi
  \AlgComment{#1}}

\AtBeginEnvironment{algorithmic}{\lineskip0pt}

\newcommand*\Let[2]{\State #1 $\gets$ #2}
\newcommand*\Stateh{\State \algstrut[1]}
\newcommand\HState[1]{\BeginBox[fill=yellow]\State{#1}\EndBox}

%% END FIXES
%% COLOURS FOR ALGORITHMS


\makeatletter
\newcommand{\algcolor}[2]{%
  \hskip-\ALG@thistlm\colorbox{#1}{\parbox{\dimexpr\linewidth-2\fboxsep}{\hskip\ALG@thistlm\relax #2}}%
}
\newcommand{\algemph}[1]{\algcolor{GreenYellow}{#1}}
\makeatother

\definecolor{dkgreen}{rgb}{0,0.6,0}
\definecolor{gray}{rgb}{0.5,0.5,0.5}
\definecolor{mauve}{rgb}{0.58,0,0.82}
\lstdefinestyle{SQL}{%
  language=SQL,
  basicstyle={\small\ttfamily},
  belowskip=3mm,
  breakatwhitespace=true,
  breaklines=true,
  classoffset=0,
  columns=flexible,
  commentstyle=\color{dkgreen},
  framexleftmargin=0.25em,
  frameshape={}{y}{}{}, %To remove to vertical lines on left, set `frameshape={}{}{}{}`
  keywordstyle=\color{blue},
  numbers=left, %If you want line numbers, set `numbers=left`
  numberstyle=\tiny\color{gray},
  showstringspaces=false,
  stringstyle=\color{mauve},
  tabsize=3,
  xleftmargin =1em
}
\lstdefinestyle{C}{%
  language=C,
  basicstyle={\small\ttfamily},
  belowskip=3mm,
  breakatwhitespace=true,
  breaklines=true,
  classoffset=0,
  columns=flexible,
  commentstyle=\color{dkgreen},
  framexleftmargin=0.25em,
  frameshape={}{y}{}{}, %To remove to vertical lines on left, set `frameshape={}{}{}{}`
  keywordstyle=\color{blue},
  numbers=left, %If you want line numbers, set `numbers=left`
  numberstyle=\tiny\color{gray},
  showstringspaces=false,
  stringstyle=\color{mauve},
  tabsize=3,
  xleftmargin =1em
}

% Reusable, self-contained C++26 snippets.  The stock listings bundle only
% provides older C++ dialects, so derive the current dialect explicitly.  The
% separate keyword classes let contextual identifiers, attributes, and common
% fixed-width types remain visually distinct from reserved words.
\definecolor{BioseqCodeBackground}{HTML}{F6F8FA}
\definecolor{BioseqCodeBorder}{HTML}{D0D7DE}
\definecolor{BioseqCodeText}{HTML}{24292F}
\definecolor{BioseqCodeKeyword}{HTML}{8250DF}
\definecolor{BioseqCodeContext}{HTML}{0550AE}
\definecolor{BioseqCodeAttribute}{HTML}{953800}
\definecolor{BioseqCodeType}{HTML}{116329}
\definecolor{BioseqCodeString}{HTML}{0A3069}
\definecolor{BioseqCodeComment}{HTML}{57606A}
\definecolor{BioseqCodePreprocessor}{HTML}{A40E26}

\lstdefinelanguage[26]{C++}[11]{C++}{%
  morekeywords={char8_t,concept,consteval,constinit,co_await,co_return,
    co_yield,contract_assert,requires},
  morekeywords=[2]{final,import,module,override,post,pre},
  morekeywords=[3]{assume,deprecated,fallthrough,indeterminate,likely,
    maybe_unused,nodiscard,no_unique_address,noreturn,unlikely},
  morekeywords=[4]{int8_t,int16_t,int32_t,int64_t,uint8_t,uint16_t,
    uint32_t,uint64_t,int_fast8_t,int_fast16_t,int_fast32_t,int_fast64_t,
    uint_fast8_t,uint_fast16_t,uint_fast32_t,uint_fast64_t,int_least8_t,
    int_least16_t,int_least32_t,int_least64_t,uint_least8_t,uint_least16_t,
    uint_least32_t,uint_least64_t,intmax_t,uintmax_t,intptr_t,uintptr_t,
    ptrdiff_t,size_t}
}

\lstdefinestyle{bioseq-cpp26}{%
  language={[26]C++},
  basicstyle=\small\ttfamily\color{BioseqCodeText},
  keywordstyle=\bfseries\color{BioseqCodeKeyword},
  keywordstyle=[2]\bfseries\color{BioseqCodeContext},
  keywordstyle=[3]\color{BioseqCodeAttribute},
  keywordstyle=[4]\color{BioseqCodeType},
  commentstyle=\itshape\color{BioseqCodeComment},
  stringstyle=\color{BioseqCodeString},
  directivestyle=\color{BioseqCodePreprocessor},
  backgroundcolor=\color{BioseqCodeBackground},
  frame=single,
  rulecolor=\color{BioseqCodeBorder},
  framerule=.4pt,
  framesep=6pt,
  aboveskip=.75\baselineskip,
  belowskip=.75\baselineskip,
  columns=fullflexible,
  keepspaces=true,
  showstringspaces=false,
  breaklines=true,
  breakatwhitespace=true,
  tabsize=2,
  escapeinside={(*@}{@*)}
}

% Put a short explanatory note into a fixed-width column on the right of a
% listing line.  It is injected through listings' escapeinside delimiters, so
% ordinary C++ comments retain normal syntax highlighting and TeX-sensitive
% characters in unrelated snippets are not interpreted as LaTeX.
\newcommand{\BioseqCodeNote}[1]{%
  \unskip\penalty0\hspace*{\fill}%
  \makebox[.44\linewidth][l]{%
    \textcolor{BioseqCodeBorder}{\rule[-.12em]{.4pt}{1.05em}}\hspace{.6em}%
    \normalfont\footnotesize\sffamily\color{BioseqCodeComment}#1}%
}

\lstnewenvironment{BioseqCppSnippet}[1][]
  {\lstset{style=bioseq-cpp26,#1}}
  {}

\lstdefinelanguage{BioseqAArch64}{%
  sensitive=true,
  alsoletter={.},
  morekeywords={add,and,bl,eor,ldp,ldr,lsl,lsr,mov,movi,orr,ret,stp,str,
    sub,ushr,ushr.4s,movi.4s,movi.2d,and.16b},
  morekeywords=[2]={sp,lr,fp,x0,x1,x2,x3,x4,x5,x6,x7,x8,x9,x10,x11,
    x12,x13,x14,x15,x16,x17,x18,x19,x20,x21,x22,x23,x24,x25,x26,x27,
    x28,x29,x30,q0,q1,q2,q3,q4,q5,q6,q7,v0,v1,v2,v3,v4,v5,v6,v7},
  morecomment=[l]{//},
  morecomment=[l]{;}
}

\lstdefinestyle{bioseq-asm}{%
  style=bioseq-cpp26,
  language=BioseqAArch64,
  keywordstyle=\bfseries\color{BioseqCodeKeyword},
  keywordstyle=[2]\color{BioseqCodeType},
  commentstyle=\itshape\color{BioseqCodeComment}
}

\lstnewenvironment{BioseqAsmSnippet}[1][]
  {\lstset{style=bioseq-asm,#1}}
  {}

\definecolor{BioseqFastxHeader}{HTML}{0969DA}
\definecolor{BioseqFastxSeparator}{HTML}{8250DF}
\definecolor{BioseqFastxQuality}{HTML}{9A6700}
\definecolor{BioseqFastxUnknown}{HTML}{57606A}

\newcommand{\FastxHeaderLine}[1]{%
  \normalfont\ttfamily\color{BioseqFastxHeader}\detokenize{#1}}
\newcommand{\FastxSeparatorLine}[1]{%
  \normalfont\ttfamily\color{BioseqFastxSeparator}\detokenize{#1}}
\newcommand{\FastxQualityLine}[1]{%
  \normalfont\ttfamily\color{BioseqFastxQuality}\detokenize{#1}}

\lstdefinestyle{bioseq-fastx-base}{%
  basicstyle=\small\ttfamily\color{BioseqFastxQuality},
  backgroundcolor=\color{BioseqCodeBackground},
  frame=single,
  rulecolor=\color{BioseqCodeBorder},
  framerule=.4pt,
  framesep=6pt,
  aboveskip=.75\baselineskip,
  belowskip=.75\baselineskip,
  columns=fullflexible,
  keepspaces=true,
  showstringspaces=false,
  breaklines=true,
  escapeinside={(*@}{@*)},
  literate={A}{{\textcolor{Plum}{\bfseries A}}}{1}
    {C}{{\textcolor{MidnightBlue}{\bfseries C}}}{1}
    {G}{{\textcolor{OrangeRed}{\bfseries G}}}{1}
    {T}{{\textcolor{Green!70!black}{\bfseries T}}}{1}
    {N}{{\textcolor{BioseqFastxUnknown}{\bfseries N}}}{1}
}

\lstdefinestyle{bioseq-fastx-fasta}{style=bioseq-fastx-base,
  basicstyle=\small\ttfamily\color{BioseqCodeText}}
\lstdefinestyle{bioseq-fastx-fastq}{style=bioseq-fastx-base}

\lstnewenvironment{FastxFile}[1][fasta]
  {\Needspace{8\baselineskip}\lstset{style=bioseq-fastx-#1}}
  {}

\lstdefinelanguage{DNA}{
    keywords = {A,C,T,G,\$}
}
\lstdefinestyle{DNA}{%
    language=DNA,
    literate={A}{{\textcolor{Plum}{A}}}{1}
    {C}{{\textcolor{MidnightBlue}{C}}}{1}
    {T}{{\textcolor{Green}{T}}}{1}
    {G}{{\textcolor{OrangeRed}{G}}}{1}
    {\$}{{\$}}{1},
    basicstyle=\large\ttfamily,
    columns=fixed,
}

% END COLOURS


\newcommand{\randeq}[1]{% 
\pgfmathparse{(int(random(-40, 40))+100)/100 * #1}%
\pgfmathresult%
}%

% \smartdiagramset{planet font=\sffamily\Huge,
%    satellite font=\sffamily,
%    planet color=green!40, uniform connection color=true,
%    uniform color list = red!40 for 8 items}
% \tikzset{satellite/.append style={regular polygon,
%   regular polygon sides=8, inner sep=0pt},
%   planet/.append style={regular polygon,
%   regular polygon sides=8, inner sep=6pt}}

\tikzset{
    bigbox/.style={draw, rounded corners, minimum width=1.5cm, minimum height=1cm},
    smallbox/.style={draw, rounded corners, minimum width=1.25cm, minimum height=0.75cm},
    tinybox/.style={draw, rounded corners, minimum width=1.25cm, minimum height=0.6cm},
    bigcircle/.style={draw, circle, minimum size=1cm},
    bigellipse/.style={draw, ellipse, minimum width=1.5cm, minimum height=1.25cm},
    place/.style={inner sep=0pt, outer sep=0pt},
    colprimaryl/.style={draw=NavyBlue, fill=LightSkyBlue, text=black},
    colprimary/.style={draw=NavyBlue, fill=NavyBlue, text=white},
    operation/.style={draw=FireBrick, fill=LightSalmon, text=black, rounded corners},
    fork/.style={decorate, decoration={show path construction, lineto code={
        \draw[->](\tikzinputsegmentfirst)-|($(\tikzinputsegmentfirst)!.5!(\tikzinputsegmentlast)$)|-(\tikzinputsegmentlast);}
    }},
    center coordinate/.style={
        execute at end picture={
        \path ([rotate around={180:#1}]perpendicular cs: horizontal line through={#1},
                                    vertical line through={(current bounding box.east)})
                ([rotate around={180:#1}]perpendicular cs: horizontal line through={#1},
                                    vertical line through={(current bounding box.west)});}},
    place/.style={
        circle,
        thick,
        draw=blue!75,
        fill=blue!20,
        minimum size=6mm,
    },
    transitionH/.style={
        rectangle,
        thick,
        fill=black,
        minimum width=8mm,
        inner ysep=2pt
    },
    transitionV/.style={
        rectangle,
        thick,
        fill=black,
        minimum height=8mm,
        inner xsep=2pt
    }
}

\let\oldtikzpicture\tikzpicture
\let\oldendtikzpicture\endtikzpicture

\renewenvironment{tikzpicture}{%
    \if\dataeq{draft}{true}\comment%
    \else\expandafter\oldtikzpicture%   
    \fi
}{%
    \if\dataeq{draft}{true}\endcomment%
    \else\oldendtikzpicture%
    \fi
}

%% Custom definitions
\makeatletter
\newcommand{\defeq}{\vcentcolon=}
\newcommand{\@giventhatstar}[2]{\left(#1\;\middle|\;#2\right)}
\newcommand{\@giventhatnostar}[3][]{#1(#2\;#1|\;#3#1)}
\newcommand{\giventhat}{\@ifstar\@giventhatstar\@giventhatnostar}
\newcommand\iverson[1]{1_{\llbracket #1 \rrbracket}}
\newcommand\problike[2]{\def\probname{#1}\def\probp1{#2}\problikeaux}
\newcommand\problikeaux[1][]{%
  \ifx\relax#1\relax%
    \probname(\probp1)
  \else%
    \probname\giventhat{\probp1}{#1}
  \fi
}
\newcommand\prob{\problike{Pr}}
\newcommand\pre{\problike{H}}
\newcommand\pri{\problike{I}}
\newcommand\prt[2]{\problike{P}{#1 \rightarrow #2}}
\newcommand\prtf[3]{\problike{P_{#1}}{#2 \rightarrow #3}}
\DeclareRobustCommand\iff{\;\Longleftrightarrow\;}

\newcommand\bigforall{%
  \mathop{\lower0.75ex\hbox{\ensuremath{%
    \mathlarger{\mathlarger{\mathlarger{\mathlarger{\forall}}}}}}}%
  \limits}

\newcommand\bigexists{%
  \mathop{\lower0.75ex\hbox{\ensuremath{%
    \mathlarger{\mathlarger{\mathlarger{\mathlarger{\exists}}}}}}}%
  \limits}

\def\smallunderbrace#1{\mathop{\vtop{\m@th\ialign{##\crcr
    $\hfil\displaystyle{#1}\hfil$\crcr
    \noalign{\kern3\p@\nointerlineskip}%
    \tiny\upbracefill\crcr\noalign{\kern3\p@}}}}\limits}

\newcommand*\BitAnd{\mathbin{\&}}
\newcommand*\BitOr{\mathbin{|}}
\newcommand*\ShiftLeft{\ll}
\newcommand*\ShiftRight{\gg}
\newcommand*\BitNeg{\ensuremath{\mathord{\sim}}}

\newcommand\bithead[2]{\bitbox[]{#1}{\scriptsize\bfseries #2}}

\DeclareRobustCommand{\bbone}{\text{\usefont{U}{bbold}{m}{n}1}}
\DeclareMathOperator{\EX}{\mathbb{E}}% expected value


\newcommand{\pluseq}{\mathrel{+}=}
\newcommand{\asteq}{\mathrel{*}=}

\newcommand{\chapterbrief}[2] {
\item \hyperref[#1]{\textbf{Chapter~\ref*{#1}}} #2
}


\makeatother

% Meta

\hypersetup{
  colorlinks=true,
  linkcolor=blue,
  filecolor=magenta,      
  urlcolor=cyan,
  pdfinfo={%
      Title={\buildpdftitle},%
      Author={\buildpdfauthor},%
      Subject={\buildpdfsubject},%
      Producer={\buildpdfproducer},%
      Creator={\buildpdfcreator},%
      Keywords={\buildpdfkeywords},%    
    },%
    keeppdfinfo %keep info entries
}
